\documentclass{book}
\usepackage[T1]{fontenc}
\usepackage[utf8]{inputenc}
\usepackage[brazil]{babel}
\usepackage{graphicx}
\usepackage{amssymb}
\usepackage{multicol}
\usepackage{hyperref}
\usepackage[a4paper, margin=2cm]{geometry}
\usepackage[framemethod=tikz]{mdframed}
% \usepackage{framed}
\usepackage{titlesec}
\usepackage{enumitem}
\usepackage{xifthen}
\usepackage{pdflscape}

\title{Livro de receitas} \author{Karl Jan Clinckspoor} \date{\today}

% Configura o espaçamento dos itens
\setitemize{itemsep=0pt}
\setenumerate{itemsep=0pt}
% Configura o cabeçalho e rodapé
\pagestyle{empty}


% Define se é para colocar figuras e emojis na versão final
\newcommand{\incluirFigura}{true}   
% Cria função para checar se é para colocar figura
\newcommand{\condicionalFigura}[1]{
  \ifthenelse{\equal{\incluirFigura}{true}}%
  {#1}%
  {}
}

% Define como que as receitas são formatadas
\titleformat
    {\section}  % command
    % [display] % shape
    {\Large\bfseries}  % format
    % {Receita \thesection:}  % label
    {}  % label
    {0.5ex}  % sep
    {\centering}  % before-code
    [\vspace{1cm}] % after-code

% Quadros customizados
% Se for para usar figuras, colocar o fundo cinza, senão, sem fundo.
\ifthenelse{\equal{\incluirFigura}{true}}{
  \newmdenv[roundcorner=10pt, backgroundcolor=gray!10,frametitle=Ingredientes,skipabove=0pt]{Ingredientes}
}{
  \newmdenv[roundcorner=10pt, backgroundcolor=white,frametitle=Ingredientes,skipabove=0pt]{Ingredientes}
}
\newmdenv[frametitle=Preparo, skipabove=0pt, linewidth=0]{Preparo}

% Configurar os emojis
\newdimen\emojiheight
\emojiheight=1cm
\newcommand\pimenta{\includegraphics[height=\emojiheight]{./emojis/pimenta}}
\newcommand{\testado}{\includegraphics[height=\emojiheight]{./emojis/checkmark}}
\newcommand{\karlAprova}{\includegraphics[height=\emojiheight]{./emojis/karl_aprova}}
\newcommand{\karenAprova}{\includegraphics[height=\emojiheight]{./emojis/karen_aprova}}

% Configurar os comandos de receita
% Receita simples, sem os emojis
\newcommand{\receita}[3]{
  \vspace*{0.1cm}
  \section{#1}
  \begin{minipage}[t]{0.3\textwidth} %
    \begin{Ingredientes}
      \noindent
      #2
    \end{Ingredientes}
  \end{minipage}%
  \quad  \begin{minipage}[t]{0.65\textwidth}
    \noindent \begin{Preparo} #3    \end{Preparo}
  \end{minipage}
  \clearpage
}

% Receita com uma caixa adicional de emojis
\newcommand{\receitaemoji}[4][]{
  \vspace*{0.1cm}
  \section{#2}

  % Se é para imprimir figuras, coloca emoji. Senão, não.
  \condicionalFigura{
    % Se tem emoji, coloca na caixa. Senão, fica vazio
    \ifthenelse{%
      \isempty{#1}
    }%
    {}%
    {%
      \begin{mdframed}[userdefinedwidth=4cm, align=center, linewidth=0]
        \centering
        #1
      \end{mdframed}
    }
  }

  \begin{minipage}[t]{0.3\textwidth} %
    \begin{Ingredientes}
      \noindent
      #3
    \end{Ingredientes}
  \end{minipage}%
  \quad  \begin{minipage}[t]{0.65\textwidth}
    \noindent \begin{Preparo} #4    \end{Preparo}
  \end{minipage}
  \clearpage
}

% Entender pq não está funcionando
\newcommand{\receitatrescolsemoji}[6]{
  \vspace*{0.1cm}
  \section{#2}

  % Se for para colocar figura, coloca emoji
  \condicionalFigura{
    % Se tem emoji, colocar, senão, ficar em branco
    \ifthenelse{%
      \isempty{#1}
    }%
    {}%
    {%
      \begin{mdframed}[userdefinedwidth=4cm, align=center, linewidth=0]
        \centering
        #1
      \end{mdframed}
    }
  }

  % Começar os ingredientes
  \begin{minipage}[t]{0.3\textwidth} %
    \begin{Ingredientes}
      \noindent
      #3
    \end{Ingredientes}
  \end{minipage}%
  \quad \begin{minipage}[t]{0.3\textwidth} %
    \begin{Ingredientes}
      \noindent
      #4
    \end{Ingredientes}
  \end{minipage}%
  \quad \begin{minipage}[t]{0.3\textwidth} %
    \begin{Ingredientes}
      \noindent
      #5
    \end{Ingredientes}
  \end{minipage}
  
  \begin{minipage}[t]{\textwidth}
    \noindent \begin{Preparo} #6    \end{Preparo}
  \end{minipage}
  \clearpage
}

% Para receitas com muitos ingredientes, que precisam de cuidados especiais
\newcommand{\receitatrescols}[5]{
  \vspace*{0.1cm}
  \section{#1}
  \begin{minipage}[t]{0.3\textwidth} %
    \begin{Ingredientes}
      \noindent
      #2
    \end{Ingredientes}
  \end{minipage}%
  \quad \begin{minipage}[t]{0.3\textwidth} %
    \begin{Ingredientes}
      \noindent
      #3
    \end{Ingredientes}
  \end{minipage}%
 \quad \begin{minipage}[t]{0.3\textwidth} %
    \begin{Ingredientes}
      \noindent
      #4
    \end{Ingredientes}
  \end{minipage}
  \vfill
  
  \begin{minipage}[t]{\textwidth}
    \noindent \begin{Preparo} #5    \end{Preparo}
  \end{minipage}
  \clearpage
}

% Comando para incluir uma foto de receita sem ter que usar o includegraphics.
% Útil no futuro caso eu queira padronizar alguma coisa
\newcommand{\fotoreceita}[2]{
  \condicionalFigura{
  \begin{center}
    \includegraphics[width=#1]{#2}
  \end{center}
  }
}

% Comando para facilitar colocar temperaturas
\newcommand{\grau}{$^\circ$}